\section{Обобщённые координаты ячейки и механическая лестница от одной ступени действия}
\label{sec:coords-dh-dt}
\label{ch:native-coords}

Вернёмся к вопросу, оставленному открытым при обсуждении макроскопических уравнений (§\ref{sec:macro-known}):
почему в механике все наблюдаемые величины раскладываются по трём базовым размерностям --- длины, времени и массы --- и откуда берутся первичные единицы импульса, энергии и силы?
На дискретной решётке ответ становится прозрачным: как только зафиксированы элементарный шаг пространства~$\caStepSpace$, такт времени~$\htick$ и единственный квант фазового действия на ячейку $\ladderAction = \planckHbar/2$, все остальные механические кванты получаются из~$\ladderAction$ делением на~$\caStepSpace$ и~$\htick$ без введения новых сущностей.

\begin{definition}[Внешние обобщённые координаты \texorpdfstring{$\caCoordsDhDt$}{dh,dt} и лестница квантов]
\label{def:dh-dt}
Пусть $\caStepSpace$ --- шаг решётки, а $\htick$ --- длительность одного такта.
\textbf{Внешними координатами} на микроуровне~$\layerMicro$ называется пара элементарных приращений
\[
  dh := \caStepSpace,\qquad dt := \htick,
\]
отношение которых задаёт тактовую скорость переноса по связи $\caSignalSpeed = dh/dt = \caStepSpace/\htick$ (в безразмерном счёте автомата полагают $dh = dt = 1$).
Из одной \textbf{ступени действия} $\ladderAction$ на ячейке~$(\hv, 1\,\htick)$ непосредственно определяются первичные кванты импульса, энергии, силы и инерции:
\[
  \ladderMomentum = \frac{\ladderAction}{dh},\qquad
  \ladderEnergy = \frac{\ladderAction}{dt},\qquad
  \ladderForce = \frac{\ladderAction}{dh\cdot dt},\qquad
  m_{\arg} = \frac{2\ladderAction\,dt}{dh^2}.
\]
\end{definition}

\begin{definition}[Внутренние координаты регистра ячейки]
\label{def:internal-coords}
Состояние внутри самой ячейки~$\hv$ полностью задаётся тройкой целых чисел $(\symMu, k_{\symPhi}, \phaseVarphi_{\mathrm{f}})$ --- или, что равносильно, парой $(\symMu, \phaseVarphi_{\mathrm{disc}})\in\spaceZ_{64}\times\spaceZ_{\phaseRingCount}$ --- для каждой из двух спинорных компонент (\cref{sec:descent-phase-disc}).
Невозбуждённому фону кипящего вакуума в этих координатах отвечает нулевой подпороговый остаток $\phaseVarphi_{\mathrm{f}} = 0$, то есть целочисленная сетка секторов $\phaseVarphi_{\mathrm{disc}} = k_{\symPhi}\cdot\symCapitalDelta\symPhi_{\mathrm{disc}} = 41\,k_{\symPhi}$.
\end{definition}

\begin{table}[ht]
\centering
\caption{Механическая лестница квантов от одной ступени действия $\ladderAction$ в координатах $\caCoordsDhDt$.}
\label{tab:dh-dt-bridge}
\begin{tabular}{@{}lll@{}}
\toprule
Величина на $\layerMicro$ & Запись через $\caCoordsDhDt$ & Выражение через $\planckHbar$ на $\layerMacro$ \\
\midrule
квант действия $\ladderAction$ & на $(\hv,1\,\htick)$ & $\planckHbar/2$ \\
квант импульса $\ladderMomentum$ & $\ladderAction/dh$ & $\planckHbar/(2\caStepSpace)$ \\
квант энергии $\ladderEnergy$ & $\ladderAction/dt$ & $\planckHbar/(2\htick)$ \\
квант силы $\ladderForce$ & $\ladderAction/(dh\cdot dt)$ & $\planckHbar/(2\caStepSpace\,\htick)$ \\
тактовая скорость $\caSignalSpeed$ & $dh/dt$ & $\caStepSpace/\htick$ \\
\bottomrule
\end{tabular}
\end{table}
